\section{Why positive iteration does not immediately remove the computation split}

Consider the staged numerical term language with variables, $0$, successor,
and the positive iterator
\[
 R^+(n,a;v.h)=h^{n+1}(a).
\]
All passive parameters of $h$ are fixed during iteration.  The defining
equations are
\[
 R^+(0,a;v.h)=h[a/v],\qquad
 R^+(Sn,a;v.h)=h[R^+(n,a;v.h)/v].
\]
This is not the tied-seed language considered in the existing drafts:
$R^+(n,a;v.Sv)=n+a+1$ has unbounded dependence on both $n$ and $a$.

\paragraph{Zero-profile theorem.}
For each fixed term $t(x_1,\ldots,x_r)$, its zero set in the standard
natural-number interpretation is exactly one of the following:
\[
 \varnothing,\qquad \mathbb N^r,\qquad
 \{\vec x:x_j=0\}\quad\text{for a single }j\in\{1,\ldots,r\}.
\]
The profile and its coordinate, when present, are effectively obtained
by structural recursion on the finite term.

\paragraph{Proof.}
Zero is always zero, a successor term is never zero, and a variable $x_j$
has zero set $x_j=0$.  Consider an iterator whose update is $h(v,\vec x)$.
By structural induction its update has one of four kinds of zero profile.
If the update is never zero, then the iterator is never zero.  If the
update is always zero, then the iterator is always zero.  If the update
is zero exactly when a single passive variable $x_j$ is zero, then the
iterator has that same zero profile.  Each of these three cases uses the
fact that at least one update occurs.  Finally, if the update is zero
exactly when its state variable $v$ is zero, each update preserves the
zero status of the current state.  The final state therefore has exactly
the seed's zero profile, already classified by structural induction.
The count's zero profile is not used: the number of updates is always
positive.  This completes the external structural induction, including
arbitrarily nested update bodies. \hfill$\square$

\paragraph{Consequences.}
Every unary term has zero set equal to $\varnothing$, $\{0\}$, or all
of $\mathbb N$.  Consequently no such term defines predecessor, whose
zero set is $\{0,1\}$, nor the zero-test that is $1$ at zero and $0$
at every positive number.  An unrestricted initialized iterator cannot
be compiled by terms: its constant-update special case is the selector
$I(0,a;v.b)=a$, $I(Sn,a;v.b)=b$, and substituting $a=1,b=0$ would
give the forbidden zero-test.
Ordinary binary addition is also not term-definable: its zero set is
$\{(x,y):x=0\text{ and }y=0\}$, which is neither a single coordinate
hyperplane nor empty nor the whole domain.  The definable shifted addition
$n+a+1$ does not contradict this conclusion, because it is never zero.

\paragraph{Exact scope.}
This is a standard-model term-definability obstruction, not an
impossibility theorem for arbitrary relational interpretations or
changed domains.  In ordinary intuitionistic first-order logic with
successor coverage, injectivity, and separation $0\ne Sm$, the formula
\[
 (n=0\wedge z=a)\ \vee\
 \exists m\,(n=Sm\wedge z=R^+(m,a;v.h))
\]
defines a total unique graph for the initialized iteration of each
fixed old-language update $h$.  This observation by itself does not
give a bi-interpretation with iPRA.  One must still prove induction for
the graph-defined expanded language and closure when newly defined
updates are nested inside iteration.  In the inherited Eq/Q/D native
calculus, even the ordinary logical connectives in this graph formula
are reconstructed using a selector, so assuming them at that point
would additionally be circular.

Counter variables that also occur as passive variables in $h$, or occur
in the seed $a$, remain present in those respective profiles; only their
separate role as the iteration count is irrelevant to the zero-profile
step.

\paragraph{Implementation motivation.}
A proposed raw encoding gives a two-field numerical array $[a,h]$
the value $a$ and a three-field array $[n,a,h]$ the value
$R^+(n,a;h)$.  For a zero-or-successor counter the computation can then
be written with a uniform predecessor expansion rather than a
zero/successor conditional.  The zero-profile theorem explains why
this apparent source reduction cannot simply inherit the existing
term compiler or selector-based logical bootstrap.
